% SQUAD-CSHARP - Rapport Technique de Projet
% Compilé avec pdflatex

\documentclass[12pt,a4paper]{report}

% ============================================================
% Packages
% ============================================================
\usepackage[utf8]{inputenc}
\usepackage[french]{babel}
\usepackage[T1]{fontenc}
\usepackage{geometry}
\usepackage{graphicx}
\usepackage[table]{xcolor}  % Option table pour \rowcolor
\usepackage{fancyhdr}
\usepackage{titlesec}
\usepackage{tocloft}
\usepackage{hyperref}
\usepackage{listings}
\usepackage{enumitem}
\usepackage{tabularx}
\usepackage{booktabs}
\usepackage{float}
\usepackage{multirow}
\usepackage{array}
\usepackage{tikz}
\usepackage{amsmath}
\usepackage{amssymb}
\usepackage{microtype}  % Meilleure justification
\usepackage{url}         % Gestion URLs

% ============================================================
% Configuration de la page
% ============================================================
\geometry{
    top=2.5cm,
    bottom=2.5cm,
    left=3cm,
    right=2.5cm,
    headheight=15pt
}

% Anti-débordement
\tolerance=1000
\emergencystretch=3em
\hfuzz=0.5pt
\sloppy

% ============================================================
% Couleurs
% ============================================================
\definecolor{primarycolor}{RGB}{34,139,34}     % Vert
\definecolor{secondarycolor}{RGB}{0,120,215}   % Bleu
\definecolor{codebackground}{RGB}{245,245,245}
\definecolor{commentcolor}{RGB}{0,128,0}
\definecolor{keywordcolor}{RGB}{0,0,255}
\definecolor{stringcolor}{RGB}{163,21,21}

% ============================================================
% Configuration hyperref
% ============================================================
\hypersetup{
    colorlinks=true,
    linkcolor=primarycolor,
    filecolor=primarycolor,
    urlcolor=secondarycolor,
    citecolor=primarycolor,
    pdftitle={SQUAD-CSHARP - Rapport Technique},
    pdfauthor={Projet SQUAD-CSHARP},
    pdfsubject={Assistant Conversationnel Expert en C\#},
    pdfkeywords={RAG, C\#, Ollama, LangChain, ChromaDB}
}

% ============================================================
% Style des titres
% ============================================================
\titleformat{\chapter}[display]
    {\normalfont\huge\bfseries\color{black}}
    {\chaptertitlename\ \thechapter}{10pt}{\Huge}
\titlespacing*{\chapter}{0pt}{20pt}{20pt}
\titleformat{\section}
    {\normalfont\Large\bfseries\color{primarycolor}}
    {\thesection}{1em}{}
\titleformat{\subsection}
    {\normalfont\large\bfseries\color{secondarycolor}}
    {\thesubsection}{1em}{}

% ============================================================
% En-têtes et pieds de page
% ============================================================
\pagestyle{fancy}
\fancyhf{}
\fancyhead[L]{\leftmark}
\fancyhead[R]{\textcolor{primarycolor}{SQUAD-CSHARP}}
\fancyfoot[C]{\thepage}
\renewcommand{\headrulewidth}{0.5pt}
\renewcommand{\footrulewidth}{0.5pt}

% Style pour la première page des chapitres
\fancypagestyle{plain}{
    \fancyhf{}
    \fancyfoot[C]{\thepage}
    \renewcommand{\headrulewidth}{0pt}
    \renewcommand{\footrulewidth}{0.5pt}
}

% ============================================================
% Configuration des listings (code)
% ============================================================
\lstset{
    backgroundcolor=\color{codebackground},
    basicstyle=\ttfamily\small,
    breaklines=true,
    captionpos=b,
    commentstyle=\color{commentcolor},
    keywordstyle=\color{keywordcolor}\bfseries,
    stringstyle=\color{stringcolor},
    numbers=left,
    numberstyle=\tiny\color{gray},
    stepnumber=1,
    numbersep=8pt,
    showstringspaces=false,
    tabsize=4,
    frame=single,
    rulecolor=\color{gray!30},
    xleftmargin=15pt,
    xrightmargin=5pt,
    framexleftmargin=10pt
}

% Styles pour différents langages
\lstdefinestyle{python}{
    language=Python,
    morekeywords={self,True,False,None,yield,async,await}
}

\lstdefinestyle{powershell}{
    language=bash,
    morekeywords={if,else,foreach,function,param}
}

% ============================================================
% Commandes personnalisées
% ============================================================
\newcommand{\checkmark}{\textcolor{green!70!black}{\checkmark}}
\newcommand{\feature}[1]{\textcolor{green!70!black}{$\checkmark$} \textbf{#1}}

% ============================================================
% Début du document
% ============================================================
\begin{document}

% ============================================================
% Page de garde
% ============================================================
\begin{titlepage}
    \centering
    
    % Informations université en haut
    \vspace*{0.5cm}
    {\large Université / École}\\[0.2cm]
    {\normalsize Formation / Filière}\\[0.2cm]
    {\normalsize Année Académique 2025-2026}\\[1cm]
    
    % Logo (à remplacer par votre image)
    \begin{center}
        \fbox{\begin{minipage}{6cm}
            \centering
            \vspace{2cm}
            \textcolor{gray}{\Large [LOGO]}\\
            \vspace{0.5cm}
            \textcolor{gray}{\small (Placer votre logo ici)}
            \vspace{2cm}
        \end{minipage}}
    \end{center}
    
    \vspace{1cm}
    
    % Titre principal
    {\huge\bfseries\color{primarycolor} SQUAD-CSHARP}\\[0.5cm]
    {\Large Rapport Technique de Projet}\\[0.3cm]
    {\large Assistant Conversationnel Expert en C\# basé sur RAG}
    
    \vspace{1.5cm}
    
    % Ligne de séparation
    \rule{\textwidth}{1.5pt}
    
    \vspace{1cm}
    
    % Informations du projet
    \begin{tabular}{ll}
        \textbf{Projet :} & SQUAD-CSHARP \\[0.3cm]
        \textbf{Version :} & 1.0 \\[0.3cm]
        \textbf{Date :} & 12 février 2026 \\[0.3cm]
        \textbf{Technologies :} & RAG, Ollama, LangChain, ChromaDB, Streamlit
    \end{tabular}
    
    \vspace{1.5cm}
    
    % Tableau des membres du groupe
    \begin{center}
        \renewcommand{\arraystretch}{1.5}
        \begin{tabularx}{0.8\textwidth}{|X|X|}
            \hline
            \rowcolor{green!20}
            \textbf{Nom \& Prénom} & \textbf{Rôle} \\
            \hline
            \textcolor{gray}{[Membre 1]} & \textcolor{gray}{[Rôle]} \\
            \hline
            \textcolor{gray}{[Membre 2]} & \textcolor{gray}{[Rôle]} \\
            \hline
            \textcolor{gray}{[Membre 3]} & \textcolor{gray}{[Rôle]} \\
            \hline
            \textcolor{gray}{[Membre 4]} & \textcolor{gray}{[Rôle]} \\
            \hline
        \end{tabularx}
    \end{center}
    
    \vfill
    

\end{titlepage}

% ============================================================
% Table des matières
% ============================================================
\tableofcontents
\newpage

% ============================================================
% CHAPITRE 1 : Introduction
% ============================================================
\chapter{Introduction}

\section{Objectif du Projet}

SQUAD-CSHARP est un chatbot intelligent spécialisé en C\# utilisant RAG (\textit{Retrieval-Augmented Generation}) pour répondre aux questions depuis une documentation locale.

\section{Contraintes Techniques}

\begin{itemize}
    \item \textbf{100\% local} sans internet
    \item \textbf{Open-source} uniquement
    \item \textbf{CPU standard}, 8 Go RAM minimum
    \item \textbf{Spécialisé C\#/.NET} uniquement
\end{itemize}

\section{Architecture Générale}

\begin{figure}[H]
    \centering
    \begin{tikzpicture}[
        node distance=1.2cm,
        box/.style={rectangle, draw, fill=green!10, text width=6cm, text centered, rounded corners, minimum height=1cm},
        arrow/.style={->, >=stealth, thick, color=primarycolor}
    ]
        \node[box] (ui) {Interface Utilisateur (Streamlit / CLI)};
        \node[box, below of=ui] (detection) {Détection \& Filtrage (prompt\_builder.py)};
        \node[box, below of=detection, xshift=-3cm] (refus) {Hors-sujet $\rightarrow$ Refus};
        \node[box, below of=detection, xshift=3cm] (conv) {Conversationnel $\rightarrow$ Direct};
        \node[box, below of=detection, yshift=-1.5cm] (rag) {Recherche RAG (retriever.py)};
        \node[box, below of=rag] (chroma) {ChromaDB (Vector Store)};
        \node[box, below of=chroma] (ollama) {Génération (Ollama + SQUAD-CSHARP)};
        \node[box, below of=ollama] (result) {Réponse Streaming + Sources};
        
        \draw[arrow] (ui) -- (detection);
        \draw[arrow] (detection) -| (refus);
        \draw[arrow] (detection) -| (conv);
        \draw[arrow] (detection) -- (rag);
        \draw[arrow] (rag) -- (chroma);
        \draw[arrow] (chroma) -- (ollama);
        \draw[arrow] (ollama) -- (result);
    \end{tikzpicture}
    \caption{Architecture du système SQUAD-CSHARP}
\end{figure}

% ============================================================
% CHAPITRE 2 : Fonctionnalités Implémentées
% ============================================================
\chapter{Fonctionnalités Implémentées}

\section{Ingestion Documentaire}

\subsection{Extraction PDF}

Utilise \texttt{PyPDF2} : lecture multi-page, extraction de texte, nettoyage automatique, conservation des métadonnées (page, source).

\subsection{Découpage Intelligent}

Chunks de 1000 caractères avec overlap de 200, séparateurs sémantiques (paragraphes, lignes, phrases).

\subsubsection{Détection Automatique du Type de Contenu}

Chaque chunk est classifié automatiquement :

\begin{table}[H]
    \centering
    \begin{tabularx}{\textwidth}{|l|X|}
        \hline
        \rowcolor{green!20}
        \textbf{Type} & \textbf{Critères de Détection} \\
        \hline
        \texttt{CODE} & Présence de patterns C\# : \texttt{class}, \texttt{namespace}, \texttt{using}, \texttt{\{}, etc. \\
        \hline
        \texttt{DEFINITION} & Patterns définitoires : "est un(e)", "désigne", "consiste à", etc. \\
        \hline
        \texttt{TEXT} & Texte descriptif standard par défaut \\
        \hline
    \end{tabularx}
    \caption{Classification automatique des types de contenu}
\end{table}

\subsection{Vectorisation et Indexation}

Modèle \texttt{all-MiniLM-L6-v2} (384 dim), ChromaDB \texttt{csharp\_docs}, métadonnées : source, page, type, timestamp.

\subsection{Statistiques d'Ingestion}

\begin{lstlisting}[style=python, caption=Exemple de sortie]
=== Statistiques d'ingestion ===
Documents traites : 13
Chunks generes : 3542
Temps d'execution : 127.3 secondes
Collection : csharp_docs
\end{lstlisting}

\section{Système RAG (Retrieval-Augmented Generation)}

\subsection{Retriever}

Recherche par similarité cosinus, top-k configurable (1-10, défaut : 5), calcul des distances, filtrage optionnel.

\subsection{Prompt Builder (Construction des Prompts)}

\subsubsection{Détection Multi-Niveau}

Le système analyse chaque message et le catégorise :

\begin{table}[H]
    \centering
    \renewcommand{\arraystretch}{1.3}
    \begin{tabularx}{\textwidth}{|l|X|l|}
        \hline
        \rowcolor{green!20}
        \textbf{Type} & \textbf{Exemples} & \textbf{RAG} \\
        \hline
        Conversationnel & Salutations, remerciements & Non \\
        \hline
        Technique C\# & Questions sur classes, LINQ, async/await & Oui \\
        \hline
        Hors-sujet & Java, Python, cuisine, histoire & Non \\
        \hline
    \end{tabularx}
    \caption{Classification des messages utilisateur}
\end{table}

\subsubsection{Gestion du Contexte}

Intégration des top-k chunks, historique des 10 derniers messages, adaptation du prompt selon l'historique.

\subsubsection{Templates Spécialisés}

3 templates : technique (avec RAG), conversationnel (minimaliste), refus (hors-sujet).

\subsection{LLM Handler}

Interface Ollama local (API REST port 11434), génération standard/streaming, gestion d'erreurs, test de connexion.

\section{Interfaces Utilisateur}

\subsection{Interface Web (Streamlit)}

Chat interactif streaming, sidebar configuration, multi-conversations, statistiques, icônes Material Symbols.

\subsubsection{Gestion des Conversations}

Création (UUID), sauvegarde JSON automatique, navigation, basculement, suppression.

\subsection{Interface CLI (Terminal)}

Version minimaliste : commandes \texttt{/reset}, \texttt{/stats}, \texttt{/exit}, coloration Rich, affichage compact.

\section{Personnalisation du Modèle}

\subsection{Modelfile SQUAD-CSHARP}

Configuration Ollama personnalisée avec :

\begin{table}[H]
    \centering
    \begin{tabular}{|l|l|l|}
        \hline
        \rowcolor{green!20}
        \textbf{Paramètre} & \textbf{Valeur} & \textbf{Effet} \\
        \hline
        \texttt{temperature} & 0.3 & Réponses précises et cohérentes \\
        \hline
        \texttt{top\_p} & 0.9 & Échantillonnage nucléaire \\
        \hline
        \texttt{top\_k} & 40 & Limitation du vocabulaire \\
        \hline
        \texttt{num\_ctx} & 4096 & Contexte étendu (tokens) \\
        \hline
    \end{tabular}
    \caption{Paramètres du modèle SQUAD-CSHARP}
\end{table}

\subsection{Règles de Comportement}

Présentation au début uniquement, refus poli hors C\#, ton pédagogique, citations des sources.

\section{Filtrage et Sécurité}

\subsection{Détection Hors-Sujet}

3 listes : mots-clés C\# (20+), autres langages (30+), génériques techniques.

\subsubsection{Logique de Priorité}

\begin{lstlisting}[style=python, caption=Algorithme de détection]
if mentions_csharp(question):
    return TECHNICAL_CSHARP
elif mentions_other_language(question):
    return OFF_TOPIC
elif has_generic_tech_keywords(question):
    return TECHNICAL_GENERIC
else:
    return CONVERSATIONAL
\end{lstlisting}

\subsection{Validation des Questions C\#}

Vérification croisée, court-circuit RAG si hors-sujet, refus strict sans infos techniques.

\section{Utilitaires}

\textbf{setup\_ollama.ps1} : Installation auto Ollama (multi-chemins, registry).

\textbf{ingest\_pdfs.py} : Ingestion avec mode reset.

\textbf{test\_connection.py} : Diagnostic complet (Ollama, ChromaDB, embeddings).

\textbf{Config} : Fichier \texttt{.env}, classe \texttt{Config}, validation au démarrage.

\textbf{Logging} : Console (INFO+), fichier \texttt{logs/squad\_csharp.log} (DEBUG+), rotation auto.

\chapter{Dépendances et Modules}

\section{Environnement Python}

Python 3.11 requis (3.14 incompatible avec ChromaDB/Pydantic v1).

\section{Dépendances Principales}

\textbf{LLM :} langchain==0.3.13, langchain-ollama==0.2.1 (interface Ollama, streaming)

\textbf{Text Splitting :} langchain-text-splitters>=0.2.0 (RecursiveCharacterTextSplitter)

\textbf{Embeddings :} sentence-transformers==3.3.1 (all-MiniLM-L6-v2, 384 dim, CPU)

\textbf{Vector DB :} chromadb==0.5.23 (stockage persistant, similarité cosinus)

\textbf{UI Web :} streamlit==1.41.1 (chat, sidebar, session\_state)

\textbf{UI CLI :} rich==13.9.4 (coloration, tables)

\textbf{PDF :} PyPDF2==3.0.1 (extraction texte)

\textbf{Config :} python-dotenv==1.0.1 (variables .env)

\section{Logiciels Externes}

\textbf{Ollama} 0.1.0+ : Serveur local LLM, port 11434

\textbf{llama3.2:3b} : Modèle base ($\sim$2 Go, 3 milliards paramètres, Meta)

% ============================================================
% CHAPITRE 4 : Difficultés Rencontrées
% ============================================================
\chapter{Difficultés Rencontrées}

\section{Python 3.14 / ChromaDB}

\textbf{Problème :} Python 3.14 incompatible avec ChromaDB (Pydantic v1).

\textbf{Solution :} Recréation venv avec Python 3.11. Temps : 30 min.

\section{Import LangChain}

\textbf{Problème :} \texttt{ModuleNotFoundError: langchain.text\_splitter}

\textbf{Solution :} Ajout \texttt{langchain-text-splitters>=0.2.0}, import depuis \texttt{langchain\_text\_splitters}. Temps : 10 min.

\section{Ollama PATH Windows}

\textbf{Problème :} \texttt{ollama} non reconnu (installé dans \texttt{\%LOCALAPPDATA\%\textbackslash Programs}).

\textbf{Solution :} Script multi-chemins + registry. Temps : 20 min.

\section{HTTP Client Closed (Streamlit)}

\textbf{Problème :} \texttt{@st.cache\_resource} sur LLMHandler fermait le client httpx après première requête.

\textbf{Solution :} Suppression du cache, création nouvelle instance à chaque requête. Temps : 45 min.

\section{Modèle Répète Présentation}

\textbf{Problème :} "Bonjour ! Je suis SQUAD-CSHARP..." à chaque réponse.

\textbf{Solution :} System prompt (« Ne te présente que lors du PREMIER échange ») + templates conditionnels selon historique. Temps : 30 min + 3 tests.

\section{Réponses Hors C\#}

\textbf{Problème :} Répondait aux questions Java/Python.

\textbf{Solution :} 3 listes (C\#, autres langages, générique), fonction \texttt{is\_off\_topic()}, filtrage avant RAG, correction regex C++ (\texttt{\textbackslash bc\textbackslash +\textbackslash +} au lieu de \texttt{\textbackslash bc\textbackslash +\textbackslash +\textbackslash b}). Temps : 1h.

\section{Conversations Non Persistantes}

\textbf{Problème :} \texttt{st.session\_state} non sauvegardé.

\textbf{Solution :} Système JSON (\texttt{data/conversations/}), fonctions \texttt{load/save/sync}, UI sidebar multi-conversations. Temps : 1h30.

\section{Émojis Interface}

\textbf{Problème :} Demande de remplacement émojis par vraies icônes.

\textbf{Solution :} Google Material Symbols (CDN), fonction helper, icônes Streamlit natives. Temps : 30 min.

\chapter{Métriques}

\section{Code}

\begin{table}[H]
    \centering
    \begin{tabular}{|l|c|c|}
        \hline
        \rowcolor{green!20}
        \textbf{Catégorie} & \textbf{Fichiers} & \textbf{Lignes} \\
        \hline
        Source (src/) & 11 & $\sim$1800 \\
        \hline
        Scripts & 3 & $\sim$400 \\
        \hline
        Tests & 2 & $\sim$300 \\
        \hline
        Config & 4 & $\sim$100 \\
        \hline
        \rowcolor{green!10}
        \textbf{TOTAL} & \textbf{20} & $\sim$\textbf{2600} \\
        \hline
    \end{tabular}
\end{table}

\section{Performance}

\textbf{Ingestion (13 PDF)} : $\sim$2 min, $\sim$3500 chunks, $\sim$150 Mo ChromaDB

\textbf{Requête RAG} : 50-150 ms (vectorielle) + 3-5 sec (LLM) = 3-5 sec total

\textbf{Mémoire} : Ollama + modèle ($\sim$2.5 Go) + ChromaDB ($\sim$200 Mo) + Streamlit ($\sim$150 Mo) = $\sim$3 Go total

\chapter{Tests et Validation}

\textbf{Tests unitaires} : \texttt{test\_ingestion.py} (extraction, découpage, vectorisation), \texttt{test\_rag.py} (recherche, prompts, détection, LLM).

\textbf{Scénarios} : Conversation simple, questions C\# avec sources, hors-sujet (refus), multi-tours, gestion conversations, réindexation.

\textbf{Cas limites} : Base vide, Ollama arrêté, modèle manquant, PDF corrompu, question vide, message très long.

\chapter{Conclusion}

\section{Objectifs Atteints}

Système RAG fonctionnel, interface moderne, filtrage C\#-only, conversations persistantes, performance standard, documentation complète.

\section{Améliorations Futures}

\textbf{Court terme :} Export conversations (PDF, MD), mode sombre, recherche historique.

\textbf{Moyen terme :} Modèle 8b + GPU, diagrammes UML, feedback utilisateur.

\textbf{Long terme :} Extension VS Code, multi-utilisateurs, fine-tuning, agent autonome.

\section{Synthèse}

SQUAD-CSHARP démontre qu'un assistant pédagogique de qualité peut être créé 100\% local avec des technologies open-source, combinant RAG, LLM local, interface moderne, filtrage intelligent et persistance, tout en respectant autonomie, coût zéro, performance standard et confidentialité.

% ============================================================
% Annexes
% ============================================================
\appendix

\chapter{Structure du Projet}

\section{Arborescence Complète}

\begin{lstlisting}[basicstyle=\ttfamily\footnotesize, numbers=none]
CHATBOT_C#/
|-- data/
|   |-- chroma_db/              # Base ChromaDB
|   |-- conversations/          # Historiques JSON
|   |-- pdfs/                   # Documents sources
|-- docs/
|   |-- SPEC_PROJECT_CSHARP_LOCAL.md
|   |-- RAPPORT_PROJET.md
|   |-- RAPPORT_TECHNIQUE.md
|-- logs/
|   |-- squad_csharp.log        # Logs applicatifs
|-- models/
|   |-- Modelfile               # Config Ollama
|-- scripts/
|   |-- ingest_pdfs.py
|   |-- setup_ollama.ps1
|   |-- test_connection.py
|-- src/
|   |-- ingestion/
|   |   |-- pdf_loader.py
|   |   |-- text_splitter.py
|   |   |-- vector_store.py
|   |-- rag/
|   |   |-- llm_handler.py
|   |   |-- prompt_builder.py
|   |   |-- retriever.py
|   |-- ui/
|   |   |-- cli_app.py
|   |   |-- streamlit_app.py
|   |-- utils/
|   |   |-- config.py
|   |   |-- logger.py
|   |   |-- validators.py
|   |-- main.py
|-- tests/
|   |-- test_ingestion.py
|   |-- test_rag.py
|-- .env.example
|-- .gitignore
|-- README.md
|-- requirements.txt
\end{lstlisting}

\chapter{Configuration Matérielle Recommandée}

\section{Configuration Minimum}

Pour tests et développement uniquement :

\begin{table}[H]
    \centering
    \renewcommand{\arraystretch}{1.3}
    \begin{tabular}{|l|l|}
        \hline
        \rowcolor{green!20}
        \textbf{Composant} & \textbf{Spécification} \\
        \hline
        CPU & 4 cœurs (Intel i5 gen 8+ ou équivalent) \\
        \hline
        RAM & 8 Go \\
        \hline
        Stockage & 10 Go disponibles \\
        \hline
        OS & Windows 10/11, Linux, macOS \\
        \hline
    \end{tabular}
    \caption{Configuration matérielle minimum}
\end{table}

\section{Configuration Recommandée}

Pour usage régulier et confortable :

\begin{table}[H]
    \centering
    \renewcommand{\arraystretch}{1.3}
    \begin{tabular}{|l|l|}
        \hline
        \rowcolor{green!20}
        \textbf{Composant} & \textbf{Spécification} \\
        \hline
        CPU & 6+ cœurs (Intel i7 / Ryzen 5) \\
        \hline
        RAM & 16 Go \\
        \hline
        Stockage & 20 Go SSD disponibles \\
        \hline
        OS & Windows 11 ou Linux récent \\
        \hline
    \end{tabular}
    \caption{Configuration matérielle recommandée}
\end{table}

\section{Configuration Optimale}

Pour performance maximale :

\begin{table}[H]
    \centering
    \renewcommand{\arraystretch}{1.3}
    \begin{tabular}{|l|l|}
        \hline
        \rowcolor{green!20}
        \textbf{Composant} & \textbf{Spécification} \\
        \hline
        CPU & 8+ cœurs (Intel i9 / Ryzen 7+) \\
        \hline
        RAM & 32 Go \\
        \hline
        GPU & NVIDIA RTX 2060+ (pour modèles plus grands) \\
        \hline
        Stockage & 50 Go SSD NVMe \\
        \hline
        OS & Windows 11 ou Ubuntu 22.04+ \\
        \hline
    \end{tabular}
    \caption{Configuration matérielle optimale}
\end{table}

\chapter{Glossaire}

\begin{table}[H]
    \centering
    \renewcommand{\arraystretch}{1.4}
    \begin{tabularx}{\textwidth}{|l|X|}
        \hline
        \rowcolor{green!20}
        \textbf{Terme} & \textbf{Définition} \\
        \hline
        \textbf{RAG} & Retrieval-Augmented Generation. Technique combinant recherche documentaire et génération de texte IA \\
        \hline
        \textbf{Embedding} & Transformation de texte en vecteur numérique permettant des calculs mathématiques \\
        \hline
        \textbf{Vector Store} & Base de données spécialisée pour stocker et chercher des vecteurs par similarité \\
        \hline
        \textbf{LLM} & Large Language Model. Modèle de langage entraîné sur énormément de texte \\
        \hline
        \textbf{Prompt} & Instruction donnée à l'IA pour qu'elle génère une réponse \\
        \hline
        \textbf{Token} & Morceau de texte (mot/sous-mot) que l'IA traite \\
        \hline
        \textbf{Streaming} & Affichage progressif de la réponse (mot par mot) pour améliorer l'UX \\
        \hline
        \textbf{Chunk} & Morceau de document découpé pour traitement et indexation \\
        \hline
        \textbf{Temperature} & Paramètre contrôlant la créativité de l'IA (0 = déterministe, 1+ = créatif) \\
        \hline
        \textbf{Top\_k} & Nombre de meilleurs résultats à conserver lors d'une recherche \\
        \hline
        \textbf{Cosine Similarity} & Mesure mathématique de similarité entre deux vecteurs (angle entre vecteurs) \\
        \hline
        \textbf{Ollama} & Serveur local permettant d'exécuter des modèles de langage (alternative à OpenAI) \\
        \hline
        \textbf{ChromaDB} & Base de données vectorielle open-source pour applications RAG \\
        \hline
    \end{tabularx}
    \caption{Glossaire des termes techniques}
\end{table}

% ============================================================
% Références
% ============================================================
\chapter*{Références}
\addcontentsline{toc}{chapter}{Références}

\begin{itemize}
    \item \textbf{Ollama} : \url{https://ollama.ai/}
    \item \textbf{LangChain} : \url{https://python.langchain.com/}
    \item \textbf{ChromaDB} : \url{https://www.trychroma.com/}
    \item \textbf{Streamlit} : \url{https://streamlit.io/}
    \item \textbf{Sentence-Transformers} : \url{https://www.sbert.net/}
    \item \textbf{Llama 3.2} : \url{https://ai.meta.com/llama/}
    \item \textbf{RAG Paper} (Lewis et al., 2020) : \url{https://arxiv.org/abs/2005.11401}
\end{itemize}

% ============================================================
% Page de fin
% ============================================================
\newpage
\thispagestyle{empty}
\vspace*{\fill}
\begin{center}
    {\Large\textcolor{primarycolor}{\textbf{SQUAD-CSHARP}}}\\[0.3cm]
    {\large Rapport Technique de Projet}\\[1cm]
    
    \rule{0.5\textwidth}{0.5pt}\\[0.5cm]
    
    {\normalsize Version 1.0 -- 12 février 2026}\\[0.3cm]
    {\normalsize $\sim$2600 lignes de code -- 20 fichiers}\\[0.3cm]
    {\normalsize Durée de développement : $\sim$15--20 heures}\\[1cm]
    
    \textit{Assistant Conversationnel Expert en C\# basé sur RAG}\\
    \textit{100\% Local -- Open-Source -- Gratuit}
\end{center}
\vspace*{\fill}

\end{document}
